\documentclass[compress,12pt]{beamer}

\usetheme{Arguelles}

\usepackage{graphicx}
\usepackage{caption}
\usepackage[spanish,es-noshorthands]{babel}
\usepackage[pages=some]{background}

\usepackage{tikz}
\usepackage{tikz-cd}
\usetikzlibrary{
  babel,
  arrows.meta,
  calc,
  positioning,
  overlay-beamer-styles,
  shadows.blur
}

\usepackage{amsmath,amssymb,latexsym,amscd} 
\usepackage[all,cmtip]{xy}
\usepackage{fancyhdr}
\usepackage{mathalfa}
\usepackage{mathrsfs}
\usepackage{hyperref}
\usepackage{ragged2e}
\usepackage{wasysym}

% Colores para la diapositiva Top/Frm
\definecolor{TopBlue}{RGB}{14,73,160}
\definecolor{FrmGreen}{RGB}{15,120,45}
\definecolor{TopLight}{RGB}{235,242,252}
\definecolor{FrmLight}{RGB}{235,248,238}

\pagenumbering{arabic}

\DeclareMathOperator{\op}{op}
\DeclareMathOperator{\pt}{pt}
\DeclareMathOperator{\spec}{spec}
\DeclareMathOperator{\Fit}{Fit}
\DeclareMathOperator{\Pth}{P}
\DeclareMathOperator{\Vrm}{Vrm}
\DeclareMathOperator{\Frm}{Frm}
\DeclareMathOperator{\Top}{Top}
\DeclareMathOperator{\Ord}{Ord}
\DeclareMathOperator{\Obj}{Obj}
\DeclareMathOperator{\Hom}{Hom}

\title{Modificaciones de parches}
\subtitle{y algunos axiomas de separación en la topología sin puntos}
%\event{58° Congreso Nacional de la SMM}
%\event{Seminario de Álgebra, CUCEI}
\date{\today}
\author{Juan Carlos Monter Cortés \\ Director: Dr. Luis Ángel Zaldívar Corichi}
%\institute{Universidad de Guadalajara}
%\email{juan.monter2902@alumnos.udg.mx}

%\homepage{www.mywebsite.com}
%\github{JCmonter}

\begin{document}


%\frame[plain]{\titlepage}

\begin{frame}[plain]
\titlepage

\vfill
\begin{center}
    \textcolor{white}{\scriptsize Investigación apoyada por SECIHTI-PROYECTO CBF2023-2024-2630}
\end{center}
\end{frame}

\begin{frame}{Contenido}
\tableofcontents
\end{frame}

\section{Resumen}

\begin{frame}{Los ingredientes principales}

\small

%--- Definición de categorías ---
\[
\begin{array}{rcl}
\Top & = &
\left\{
\begin{array}{l}
\text{Objetos: espacios topológicos} \\
\text{Flechas: funciones continuas}
\end{array}
\right. \\[0.6cm]
\Frm & = &
\left\{
\begin{array}{l}
\text{Objetos: marcos} \\
\text{Flechas: morfismos de marcos}
\end{array}
\right.
\end{array}
\]

\vspace{0.2cm}

%--- Diagrama ---
\begin{center}
\begin{tikzpicture}[node distance=4cm, auto]
\node (Top) {\Large $\Top$};
\node (Frm) [right of=Top] {\Large $\Frm$};

\draw[->] (Top) to[bend left] node[above] {$\mathcal{O}(\_)$} (Frm);
\draw[->] (Frm) to[bend left] node[below] {$\pt(\_)$} (Top);
\end{tikzpicture}
\end{center}
\end{frame}

\begin{frame}{Las herramientas}
\small
\begin{columns}[T,onlytextwidth]

\column{0.48\textwidth}
\textbf{Núcleos}
\begin{itemize}
    \item Núcleos abiertos
    \item Núcleos cerrados
    \item Núcleos regulares
    \item Núcleos ajustados
\end{itemize}

\vspace{0.4cm}

\onslide<2->{\textbf{Filtros}
\begin{itemize}
    \item Filtros abiertos
    \item Filtros completamente primos
    \item Filtros de admisibilidad
\end{itemize}}

\column{0.48\textwidth}

\onslide<3->{
\textbf{Correspondencias biyectivas}
\begin{itemize}
    \item Puntos--caracteres
    \item Núcleos--cocientes
    \item Teorema de Hofmann--Mislove
\end{itemize}}

\vspace{0.4cm}

\onslide<4->{
\textbf{Nociones topológicas}
\begin{itemize}
    \item Morfismo reflexión espacial
    \item Axiomas de separación
    \item Núcleos espacialmente inducidos
\end{itemize}}

\end{columns}
\vspace{0.4cm}

\onslide<5->{
\begin{center}
\textbf{Entre otras...}
\end{center}}
\end{frame}

\begin{frame}{Derivadas y prenúcleos}

\end{frame}

\begin{frame}[fragile]{Espacio de parches}
    Consideremos $S\in \Top$. Denotamos por $^pS=(S, \mathcal{O}^pS)$ al \textbf{espacio de parches}\footnote{El espacio de parches fue introducido por Hochster en \cite{hochster1969prime}} de $S$, donde $\mathcal{O}^pS$ está generado por
    \[
    \mbox{pbase}=\{U\cap Q'\mid U\in \mathcal{O}S, Q\in \mathcal{Q}S\}
    \]
   \[\begin{tikzcd}
	{\mathcal{O}S} & {\mathcal{O}^pS}
	\arrow[from=1-1, to=1-2]
\end{tikzcd}\]

\onslide<2->{
\begin{block}{Definición}
$S$ es \textbf{empaquetado} si todo subconjunto compacto (saturado) es cerrado.
\end{block}
\[
S \mbox{ es empaquetado}\iff\, ^pS=S\quad \mbox{ y }\quad T_2\Rightarrow \mbox{empaquetado}\Rightarrow T_1
\]}
\end{frame}

\begin{frame}[fragile]{Marco de parches}
Consideremos $A\in \Frm$. Denotamos por $PA$ al \textbf{marco de parches}\footnote{Información sobre el marco de parches puede ser consultada en \cite{sexton2006point, sexton2006ordinal}} de $A$, donde $PA$ está generado por
\[
\mbox{Pbase}=\{u_a\wedge v_F\mid a\in A, F\in A^\wedge\}
\]

\[\begin{tikzcd}
	A & PA
	\arrow[from=1-1, to=1-2]
\end{tikzcd}\]

\onslide<2->{
\begin{block}{Definición}
    $A$ es \textbf{parche trivial} si y solo si $A\cong PA$.
\end{block}
      \[
        ?\Rightarrow \mbox{ parche trivial }\Rightarrow T_1
      \]
}
\end{frame}

\begin{frame}{El diccionario}

\begin{columns}[T,onlytextwidth]
  \column{0.48\textwidth}
  \begin{block}{\centering $\Top$}
    \begin{itemize}
      \item<2-> Subespacios
      \item<4-> Subespacios cerrados y compactos
      \item<6-> Axiomas de separación 
      \item<8-> Espacio de parches ($^pS$)
      \item<10-> pbase$=\{U\cap Q'\}$ 
      \item<12-> Empaquetado ($^pS=S$)
    \end{itemize}
  \end{block}

  \column{0.48\textwidth}
  \begin{block}{\centering $\Frm$}
    \begin{itemize}
      \item<3-> Núcleos-Cocientes
      \item<5-> Cocientes cerrados y compactos
      \item<7-> Propiedades de separación
      \item<9-> Marco de parches ($PA$)
      \item<11-> Pbase$=\{u_a\wedge v_F\}$
      \item<13-> Parche trivial ($PA\cong A$)
    \end{itemize}
  \end{block}
\end{columns}
\end{frame}

\begin{frame}{Marcos eficientes}
\begin{block}{Definición \cite[Def. 8.2.1]{sexton2003point}}
Sean $A\in \Frm$, $F\in A^\wedge$ y $\alpha\in \Ord$. Decimos que:
\begin{enumerate}
    \item<2-> $F$ es $\alpha$\textbf{-eficiente} si para $x\in F$, $d\vee x=1$, donde
    \[
    d:=d(\alpha)=f^\alpha(0).
    \]
    \item<3-> $A$ es $\alpha$\textbf{-eficiente} si cada $F\in A^\wedge$ es $\alpha$-eficiente.
    \item<4-> $A$ es \textbf{eficiente} si es $\alpha$-eficiente para algún $\alpha\in \Ord$.
\end{enumerate}
\end{block}

\onslide<5->{\begin{block}{Proposición \cite[Lema 8.2.2]{sexton2003point}}
\[
    A \mbox{ es eficiente}\quad \iff\quad A \mbox{ es parche trivial}.
\]
\end{block}}
\end{frame}

\begin{frame}{Para qué sirve la eficiciencia}

\end{frame}

\begin{frame}{Propiedades de los marcos eficicientes}

\end{frame}

\begin{frame}{Eficiencia y su relación con otras construcciones}

\end{frame}

\section{Avances}

\begin{frame}{Empujes y levantamientos}
  
\end{frame}

\begin{frame}{Nociones de apilamiento}
Consideremos $S\in\Top$. Tenemos $Q\in \mathcal{Q}S$ y $F\in \mathcal{O}S^\wedge$. Al par $(F, Q)$ lo llamamos para HM.
Adicionalmente

\[
v_F\leq [Q']\leq w_F
\]

\begin{block}{Definición:}
  Sea $S\in \Top$ y $(F, Q)$ un par HM.
\begin{enumerate}
  \item Un conjunto $Q\in \mathcal{Q}S$ es apilado si $\overline{Q}=\partial_F^\infty(S)$ y es fuertemente apilado si $v_F=[Q']$.
  \item El espacio $S$ es apilado o fuertemente apilado si cada $Q\in \mathcal{Q}S$ es apilado o fuertemente apilado, respectivamente.
\end{enumerate}
\end{block}

\end{frame}

\begin{frame}{Marcos apilados y fuertemente apilados}
\begin{block}{Definición:}
Sean $A\in \Frm$, $S=\pt A$ y $(F,Q)$ un par HM. Decimos que $A$ es:
\begin{enumerate}
  \item apilado si $(\mathcal{U}_*([Q'])\mathcal{U})(0)=v_F(0)$.
  \item fuertemente apilado si $\mathcal{U}_*[Q']\mathcal{U}=v_F$.
\end{enumerate}
\end{block}

\begin{block}{Proposición:}
Las nociones de fuertemente apilado y apilado son conservativas.
\end{block}

\begin{proof}
  Verificamos que 
\[
S \text{ apilado (f. apilado)}\iff \mathcal{O}S \text{ apilado (f. apilado)}.
\]
\end{proof}
\end{frame}

\begin{frame}{Interacción de la eficiencia con otras propiedades}

\end{frame}

\begin{frame}{Lista de actividades}

\end{frame}


\End

\section*{\textsc{Referencias}}
\begin{frame}[allowframebreaks]
%\frametitle{Bibliografía}
%\begin{thebibliography}{20}\markboth{Bibliografía}{Bibliografía}

%\bibitem{P.T.} P. T. Johnstone, \textit{Stone spaces}, Cambridge Studies in Advanced Mathematics, vol. 3, Cambridge University Press, Cambridge, 1982. MR 698074

%\bibitem{J.M.} J. Monter; A. Zaldívar, \textit{El enfoque locálico de las reflexiones booleanas: un análisis en la categoría de marcos} [tesis de maestría], 2022. Universidad de Guadalajara.

%\bibitem{P.S.} J. Paseka and B. Smarda, \textit{$ T_2 $-frames and almost compact frames.} Czechoslovak Mathematical Journal (1992), 42(3), 385-402.

%\bibitem{J.P.} J. Picado and A. Pultr, \textit{Frames and locales: Topology without points}, Frontiers in Mathematics, Springer Basel, 2012.

%\bibitem{J.P.2} J. Picado and A. Pultr, \textit{Separation in point-free topology}, Springer, 2021.

%\bibitem{R.S.} RA Sexton, \textit{A point free and point-sensitive analysis of the patch assembly}, The University of Manchester (United Kingdom), 2003.

%\bibitem{R.S.2} RA Sexton, \textit{Frame theoretic assembly as a unifying construct}, The University of Manchester (United Kingdom), 2000.

%\bibitem{S.S.06} RA. Sexton and H.~Simmons, \textit{An ordinal indexed hierarchy of separation properties}, Topology Proceedings, Vol.~30, No.~2, 2006, pp.~585--625.
%\bibitem{R.S.3} RA Sexton and H. Simmons, \textit{Point-sensitive and point-free patch constructions}, Journal of Pure and Applied Algebra \textbf{207} (2006), no. 2, 433-468.

%\bibitem{H.S.} H. Simmons, \textit{An Introduction to Frame Theory}, lecture notes, University of Manchester. Disponible en línea en \url{https://web.archive.org/web/20190714073511/http://staff.cs.manchester.ac.uk/~hsimmons}.

%\bibitem{H.S.R} H. Simmons, \textit{Regularity, fitness, and the block structure of frames.} Applied Categorical Structures 14 (2006): 1-34.

%\bibitem{H.S.4} H. Simmons, \textit{The lattice theoretic part of topological separation properties}, Proceedings of the Edinburgh Mathematical Society, vol.~21, pp.~41--48, 1978.

%\bibitem{H.S.V} H. Simmons, \textit{The Vietoris modifications of a frame}. Unpublished manuscript (2004), 79pp., available online at http://www. cs. man. ac. uk/hsimmons.

%\bibitem{A.W.} A. Wilansky, \textit{Between T1 and T2}, MONTHLY (1967): 261-266.

%\bibitem{A.Z.} A. Zaldívar, \textit{Introducción a la teoría de marcos} [notas curso], 2025. Universidad de Guadalajara.
%\end{thebibliography}

\bibliographystyle{plain}
\bibliography{research2}
\end{frame}

\begin{frame}[plain,standout]
      \centering
      \Huge{\bfseries \smiley ¡Muchas gracias! \smiley}\\[1cm]
      %In combination with \textit{plain}, \\
      %it makes a nice thank-you slide!
      \scalebox{4}{\faGithub} \par\bigskip
      \url{https://github.com/JCmonter/Apuntes/tree/main/Presentaciones} \\
      %\url{https://ctan.org/pkg/beamertheme-arguelles}
\end{frame}
\end{document}